% =====================================================================
%  Inferência Bayesiana com MrBayes e delimitação de linhagens com mPTP
%  Prática 13 -- Métodos Computacionais em Sistemática e Filogeografia (UNESP)
%
%  Compilar com pdfLaTeX (Overleaf: Menu > Compiler > pdfLaTeX).
%
%  VERSÃO: por padrão gera a versão do ALUNO (sem as notas do professor).
%  Para a versão do PROFESSOR, troque \professorfalse por \professortrue
%  na linha indicada abaixo, ou compile com
%     pdflatex "\def\versaoprofessor{}\input{apostila_mrbayes_mptp.tex}"
% =====================================================================
\documentclass[11pt,a4paper]{article}

\newif\ifprofessor
\professorfalse            % <-- troque para \professortrue para incluir as notas do professor
\ifdefined\versaoprofessor \professortrue \fi

\usepackage[T1]{fontenc}
\usepackage[utf8]{inputenc}
\usepackage{lmodern}
\usepackage{textcomp}
\usepackage{amsmath}
\usepackage[provide=*,brazilian]{babel}
\usepackage[a4paper,margin=2.5cm]{geometry}
\usepackage{fancyvrb}
\usepackage{upquote}
\usepackage{booktabs}
\usepackage{array}
\usepackage{enumitem}
\usepackage{comment}
\usepackage{needspace}
\usepackage{url}
\usepackage[hidelinks]{hyperref}

\fvset{xleftmargin=1.5em,fontsize=\small}
\setlist{itemsep=0.2em,topsep=0.4em}
\setlength{\parskip}{0.35em}
\setlength{\emergencystretch}{3em}
\urlstyle{tt}

\newcommand{\arq}[1]{\path{#1}}
\newcommand{\ui}[1]{\textit{#1}}
\newcommand{\cmd}[1]{\texttt{#1}}

\ifprofessor
  \specialcomment{resposta}{\begin{quote}\small\textbf{Para o professor.}\ }{\end{quote}}
\else
  \excludecomment{resposta}
\fi

\newcounter{questao}
\newcommand{\questao}{\par\needspace{3\baselineskip}\refstepcounter{questao}\noindent\textbf{Questão~\thequestao.}\ }

\title{Inferência Bayesiana com MrBayes e delimitação de linhagens com mPTP\\[0.3em]
\large Prática 13: o ND2 de \emph{Anolis chrysolepis} e \emph{A. tandai}}
\author{Métodos Computacionais em Sistemática e Filogeografia\\Universidade Estadual Paulista (UNESP)}
\date{\ifprofessor Versão do professor (com notas)\else Versão do aluno\fi\\ Novembro de 2026}

\begin{document}
\maketitle

\paragraph{Origem do material.} A Parte~1 adapta o tutorial de MrBayes da página
\url{https://labesta.org/filogenia/bayesiana}; a Parte~2 segue a documentação do mPTP
(\url{https://github.com/Pas-Kapli/mptp}). As mudanças estão no Apêndice~\ref{ap:mudancas}.

\paragraph{Os dados.} Um subconjunto das sequências do gene mitocondrial ND2 de lagartos do grupo
\emph{Anolis chrysolepis}, publicadas por Vieira de Oliveira \emph{et al.} \cite{oliveira2025}:
110 sequências de 924 pares de bases, sendo 58 de \emph{A. chrysolepis} (nomes que começam com
\cmd{CR\_}), 50 de \emph{A. tandai} (\cmd{TA\_}) e uma de cada um dos dois grupos externos,
\emph{A. brasiliensis} (\cmd{BR\_}) e \emph{A. planiceps} (\cmd{PL\_}).

\paragraph{O roteiro.} Duas análises encadeadas, a segunda usando a saída da primeira:
\begin{enumerate}
\item \textbf{MrBayes}: a filogenia Bayesiana do ND2 (Parte~1);
\item \textbf{mPTP}: quantas linhagens a árvore sugere (Parte~2).
\end{enumerate}

\paragraph{Arquivos.} A pasta \arq{pratica_mrbayes_mptp} contém:
\begin{itemize}
\item \arq{anolis_ND2.nex}: o alinhamento em NEXUS, já com o bloco do MrBayes;
\item \arq{anolis_ND2.fas}: o mesmo alinhamento em FASTA;
\item \arq{preparar_mptp.R}: prepara a árvore do MrBayes para o mPTP;
\item \arq{anolis_mptp.tre} e \arq{anolis_mptp.fas}: a árvore e o alinhamento já preparados
pelo professor para o mPTP, a partir de uma corrida mais longa do MrBayes
(Seção~\ref{sec:preparar}).
\end{itemize}

\paragraph{Programas.} MrBayes 3.2.9 \cite{ronquist2012}, Tracer \cite{rambaut2018},
FigTree e mPTP 0.2.5 \cite{kapli2017}. Opcional: R com o pacote \cmd{ape} \cite{paradis2019},
para quem quiser preparar a própria árvore para o mPTP. As instruções de instalação estão em
cada parte.

\paragraph{Um aviso.} Esta é uma análise de tutorial, montada para rodar numa aula. Ao longo do
texto, os quadros ``\textbf{Na vida real}'' indicam o que faríamos diferente numa análise para
publicação.

\begin{resposta}
Antes da aula: (1) peça que os alunos instalem MrBayes, Tracer e FigTree;
(2) resolva o mPTP de cada um (Seção~\ref{sec:instalarmptp}): é a parte que mais dá trabalho;
(3) a corrida do MrBayes leva uns 3 minutos, mas vale explicar a analogia do andarilho antes de rodar, para os alunos saberem o que
estão vendo na tela. Os números esperados estão no gabarito.
\end{resposta}

% =====================================================================
\section{Introdução}
% =====================================================================

\subsection{Inferência Bayesiana de filogenias}

Na inferência Bayesiana, o que queremos é a probabilidade posterior de cada árvore (e de cada
valor dos parâmetros do modelo), dados os dados. Pelo teorema de Bayes,
\[
P(\text{árvore, parâmetros} \mid \text{dados}) =
\frac{P(\text{dados} \mid \text{árvore, parâmetros}) \times P(\text{árvore, parâmetros})}
{P(\text{dados})}.
\]
No numerador estão a verossimilhança (a mesma da máxima verossimilhança) e o \emph{prior}. O
denominador, a verossimilhança marginal, é a soma de verossimilhança $\times$ \emph{prior} sobre
todas as árvores e todos os valores dos parâmetros. É por causa dele que não dá para calcular o
posterior diretamente: para mais do que poucos terminais, há árvores demais. O MCMC
(\emph{Markov chain Monte Carlo}) contorna o problema porque só compara o posterior de um estado
com o de outro. Nessa razão, $P(\text{dados})$ aparece em cima e embaixo e se cancela, e basta
saber que o posterior é \emph{proporcional} ao numerador.

\paragraph{A analogia do andarilho.} Imagine todas as combinações possíveis de árvore e
parâmetros como uma paisagem de morros, em que a altura de cada ponto é o posterior. Um
andarilho começa num ponto qualquer e, a cada passo, sorteia um ponto vizinho: se o novo ponto
é mais alto, ele vai; se é mais baixo, ele vai só às vezes, com probabilidade igual à razão das
alturas (um ponto com metade da altura é aceito metade das vezes). Ele nunca precisa saber a
altura de um ponto, só quantas vezes um é mais alto que o outro, e é por isso que
$P(\text{dados})$ não faz falta. Depois de muitos passos, o tempo que o andarilho passa em cada
região é proporcional à altura dela. Anotando onde ele está de tempos em tempos, temos uma
amostra do posterior. A mesma imagem explica os três detalhes que
vão aparecer na prática:
\begin{itemize}
\item \textbf{\emph{Burn-in}.} O andarilho começa num lugar sorteado, em geral baixo, e os
primeiros passos são a subida até os morros. Essas primeiras anotações não representam o
posterior, e são descartadas (aqui, os primeiros 25\%).
\item \textbf{Cadeias aquecidas.} O risco é o andarilho ficar preso num morro e nunca cruzar o
vale até outro morro, talvez mais alto. Por isso, cada corrida tem quatro andarilhos: um
``frio'', que anda na paisagem verdadeira e é o único cujas posições anotamos, e três
``aquecidos'', que andam numa versão achatada da mesma paisagem, com vales mais rasos, e cruzam
de um morro para o outro com facilidade. De tempos em tempos, dois andarilhos propõem trocar de
lugar; assim, o frio pode pular para um morro que um aquecido encontrou. Cada andarilho é uma
\emph{cadeia}; o conjunto é o MC$^3$ (\emph{Metropolis-coupled} MCMC).
\item \textbf{Duas corridas independentes.} Uma \emph{corrida} é um grupo de quatro
andarilhos. Rodamos duas, soltas em pontos diferentes da paisagem e sem contato uma com a
outra. Se os andarilhos frios das duas acabam passando o mesmo tempo nos mesmos morros, temos
um bom sinal de que os dois exploraram a paisagem toda: as corridas \emph{convergiram}. Se cada
um passou a maior parte do tempo num morro diferente, não convergiram. É isso que vamos medir
(Seção~\ref{sec:convergencia}): o ASDSF compara as árvores das duas corridas, e o PSRF compara
os parâmetros.
\end{itemize}

No fim, a probabilidade posterior de um clado é simplesmente a fração das árvores amostradas
(depois do \emph{burn-in}) em que aquele clado aparece.

\subsection{Delimitando linhagens com uma árvore: o mPTP}

O modelo PTP (\emph{Poisson Tree Processes}) \cite{zhang2013} parte de uma ideia simples: numa
árvore de um único gene, os ramos \emph{entre} espécies e os ramos \emph{dentro} das espécies
crescem de jeitos diferentes. Os ramos entre espécies refletem a especiação, um processo lento,
com ramos longos; os ramos dentro das espécies refletem a coalescência dentro de cada população,
com ramos curtos e muitos nós. O PTP modela o número de substituições nos ramos como dois
processos exponenciais e procura o ponto da árvore em que se passa de um para o outro. Cada clado
abaixo desse ponto é uma espécie (para nós, uma \emph{linhagem}, veja abaixo).

O mPTP (\emph{multi-rate} PTP) \cite{kapli2017} permite que cada espécie tenha a sua própria
taxa do processo de dentro (a opção \cmd{-{}-multi}), já que espécies com populações de tamanhos
diferentes ou amostradas de jeitos diferentes têm ramos internos de tamanhos diferentes. O mPTP
encontra a delimitação de máxima verossimilhança (\cmd{-{}-ml}) e estima o suporte de cada nó com
um MCMC (\cmd{-{}-mcmc}).

Duas exigências do mPTP vão ditar a preparação da árvore:
\begin{itemize}
\item a árvore precisa ter comprimentos de ramo em substituições por sítio (uma árvore de
máxima verossimilhança ou Bayesiana, e não um cronograma);
\item a árvore precisa ser \emph{binária} (sem politomias) e enraizada.
\end{itemize}

\paragraph{Linhagem não é espécie.} Com um único gene mitocondrial, o mPTP delimita
\emph{linhagens mitocondriais}. Isso é uma hipótese a ser testada com outros dados (genes
nucleares, morfologia, distribuição, canto, ecologia), e não uma lista de espécies. Métodos de
um único loco tendem a dividir demais quando há estrutura populacional forte, e o DNA
mitocondrial, herdado só pela mãe e sem recombinação, conta a história de uma única molécula.

% =====================================================================
\section{Parte 1: a filogenia Bayesiana com o MrBayes}
% =====================================================================

\subsection{Instalação}

Vamos usar a versão 3.2.9 do MrBayes. Na página da versão,
\url{https://github.com/NBISweden/MrBayes/releases/tag/v3.2.9}, baixe o executável do seu
sistema.
\begin{itemize}
\item \textbf{Windows.} O arquivo é o \arq{mb_3.2.9_win_x64.exe}. \textbf{Copie esse arquivo
para a pasta do tutorial} (a pasta onde está o \arq{anolis_ND2.nex}) e abra-o de lá, com dois
cliques: o MrBayes procura os arquivos na pasta em que ele está, e não vai achar o
\arq{anolis_ND2.nex} se ficar na pasta \arq{Downloads}.
\item \textbf{Mac e Linux.} Ponha o executável na pasta do tutorial e, no Terminal, dentro
dela, deixe-o executável com \cmd{chmod +x} seguido do nome do arquivo. No Mac, rode antes
\cmd{xattr -d com.apple.quarantine} seguido do nome do arquivo, para tirar o bloqueio que o Mac
põe em programas baixados. Depois, rode-o com \cmd{./} seguido do nome (por exemplo,
\cmd{./mb}, se você renomear o arquivo para \arq{mb}).
\item \textbf{WSL.} Quem quiser usar o WSL (o Linux dentro do Windows; instalação na
Seção~\ref{sec:instalarmptp}) é só seguir as instruções do Linux, dentro do Ubuntu.
\end{itemize}
Vamos usar também o Tracer (\url{https://github.com/beast-dev/tracer/releases}) e o FigTree
(\url{https://github.com/rambaut/figtree/releases}), que precisam do Java.

\begin{resposta}
Os resultados do gabarito foram feitos com o MrBayes 3.2.9 compilado do código-fonte no Linux
(\cmd{./configure -{}-without-beagle -{}-without-mpi} e \cmd{make}), que dá o mesmo programa que os
executáveis da página. Se o executável de algum Mac não abrir, essa é a alternativa: com as
ferramentas de linha de comando (\cmd{xcode-select -{}-install}), \cmd{git clone} do repositório
\url{https://github.com/NBISweden/MrBayes}, os dois comandos acima, e o programa fica em
\arq{src/mb}.
\end{resposta}

\subsection{O arquivo NEXUS}

Abra o \arq{anolis_ND2.nex} num editor de texto simples (Bloco de Notas, TextEdit em modo texto
puro, ou o próprio RStudio). Um arquivo NEXUS é feito de blocos, cada um começando com
\cmd{begin} e terminando com \cmd{end;}. O primeiro é o bloco de dados:
\begin{Verbatim}
#NEXUS
begin data;
  dimensions ntax=110 nchar=924;
  format datatype=dna missing=? gap=-;
  matrix
CR_SP.MRT6281    AACACAATAGCAATCATCCC...
...
;
end;
\end{Verbatim}
O número de terminais (\cmd{ntax}) e de sítios (\cmd{nchar}) precisa bater com a matriz. Depois
dela vem o bloco do MrBayes, com os comandos que o programa vai executar em sequência:
\begin{Verbatim}
begin mrbayes;
  set autoclose=yes nowarn=yes;
  log start filename=anolis_ND2_log.txt replace;
  outgroup BR_ACB.CHUNB14548;

  charset ND2_pos1 = 1-924\3;
  charset ND2_pos2 = 2-924\3;
  charset ND2_pos3 = 3-924\3;
  partition codons = 3: ND2_pos1, ND2_pos2, ND2_pos3;
  set partition = codons;

  lset applyto=(all) nst=mixed rates=invgamma;
  unlink statefreq=(all) revmat=(all) shape=(all) pinvar=(all);
  prset applyto=(all) ratepr=variable;

  mcmc ngen=100000 nruns=2 nchains=4 samplefreq=100 printfreq=1000
       diagnfreq=1000 relburnin=yes burninfrac=0.25;
  sump;
  sumt;
  sumt contype=allcompat conformat=simple outputname=anolis_ND2_allcompat;
  log stop;
end;
\end{Verbatim}
Linha por linha:
\begin{itemize}
\item \cmd{set autoclose=yes nowarn=yes}: no fim da corrida, não pergunte se queremos continuar,
e não peça confirmação para sobrescrever arquivos.
\item \cmd{log start ...}: guarde tudo o que aparecer na tela no arquivo
\arq{anolis_ND2_log.txt}.
\item \cmd{outgroup}: a sequência de \emph{A. brasiliensis} enraíza as árvores que o MrBayes
escreve (o MrBayes não estima a raiz; isso é só para desenhar).
\item \cmd{charset} e \cmd{partition}: o ND2 é um gene codificador de proteína, e as três
posições do códon evoluem de jeitos muito diferentes (a terceira, em geral, muito mais rápido).
\cmd{1-924\textbackslash 3} quer dizer ``do sítio 1 ao 924, de três em três''. As três posições
viram três partições, cada uma com o seu modelo.
\item \cmd{lset nst=mixed rates=invgamma}: o modelo de substituição de cada partição. Com
\cmd{nst=mixed}, o MrBayes não usa um modelo fixo: ele \emph{amostra} entre todos os modelos da
família do GTR (do JC ao GTR) durante a própria corrida, em proporção à probabilidade posterior
de cada um \cite{huelsenbeck2004}. Assim, não precisamos escolher o modelo antes. \cmd{rates=invgamma}
inclui a variação de taxas entre sítios (gama) e uma proporção de sítios invariáveis.
\item \cmd{unlink}: por padrão, as partições compartilham os parâmetros; aqui, cada partição tem
as suas próprias frequências de bases, taxas de substituição, forma da gama e proporção de
invariáveis.
\item \cmd{prset ratepr=variable}: cada partição tem a sua taxa média de evolução (a terceira
posição pode evoluir mais rápido que as outras).
\item \cmd{mcmc}: a corrida. \cmd{ngen}: número de gerações (os passos do andarilho), aqui
100 mil, para caber na aula; \cmd{nruns}: corridas independentes; \cmd{nchains}: cadeias
(andarilhos) por corrida; \cmd{samplefreq}: amostra a cada 100 gerações (1.000 árvores por
corrida, no total); \cmd{printfreq}: mostra o andamento na tela a cada 1.000 gerações;
\cmd{diagnfreq}: calcula o diagnóstico de convergência a cada 1.000; \cmd{relburnin} e
\cmd{burninfrac}: descarta os primeiros 25\% das amostras.
\item \cmd{sump}: resume os parâmetros (arquivos \arq{.p}).
\item \cmd{sumt}: resume as árvores (arquivos \arq{.t}) numa árvore de consenso
(Seção~\ref{sec:consenso}).
\item o segundo \cmd{sumt}: um consenso diferente, que só vai servir para o mPTP
(Seção~\ref{sec:consenso}).
\end{itemize}

\questao Abra o MrBayes e digite \cmd{help}. Depois, \cmd{help lset} e \cmd{help mcmc}. Qual é o
valor padrão de \cmd{nchains}? E o que faz a opção \cmd{temp}?

\begin{resposta}
O padrão de \cmd{nchains} é 4, e o de \cmd{nruns} é 2; \cmd{temp} é o ``calor'' das cadeias
aquecidas (padrão 0,1): quanto maior, mais livres elas andam. O \cmd{help lset} mostra a tabela
com os valores atuais de cada opção, e é um bom jeito de conferir o que o bloco fez.
\end{resposta}

\subsection{Rodando}

Abra o MrBayes (no Windows, com dois cliques no executável que você copiou para a pasta do
tutorial; no Mac, no Linux e no WSL, com \cmd{./} e o nome do executável no Terminal, dentro da
pasta) e digite:
\begin{Verbatim}
execute anolis_ND2.nex
\end{Verbatim}
O MrBayes lê os dados, monta o modelo e começa a corrida. A tela mostra o andamento a cada
1.000 gerações:
\begin{Verbatim}[fontsize=\footnotesize]
 50000 -- [-7747.075] (-7790.338) (-7770.700) (-7820.862) * (-7773.493) ... -- 0:01:33
\end{Verbatim}
O primeiro número é a geração; depois vêm o logaritmo da verossimilhança das quatro cadeias da
primeira corrida, um asterisco e as quatro da segunda. Os colchetes marcam a cadeia fria (o andarilho
cujas posições anotamos); quando ela muda de lugar na linha, é porque dois andarilhos trocaram
de posição. No fim da linha, o tempo que falta. A cada 1.000 gerações, aparece também
\begin{Verbatim}
 Average standard deviation of split frequencies: 0.037924
\end{Verbatim}
que é o nosso primeiro diagnóstico de convergência (próxima seção).

A corrida de 100 mil gerações leva uns 3 minutos. Enquanto ela roda, leia a próxima seção.

\begin{resposta}
O resto da prática não depende da corrida do aluno, porque a árvore do mPTP é a do professor.
Todas as versões, como instaladas aqui, usam um núcleo só.
\end{resposta}

\subsection{A corrida convergiu?}\label{sec:convergencia}

Três diagnósticos, do mais rápido ao mais completo:
\begin{enumerate}
\item \textbf{O desvio padrão médio das frequências dos clados (ASDSF).} Para cada clado, o
MrBayes compara a sua frequência nas duas corridas. Se as duas amostram as mesmas árvores, as
frequências são parecidas, e a média dos desvios padrão tende a zero. O critério mais usado é
ASDSF $< 0{,}01$; entre 0,01 e 0,05, a corrida pode servir para uma ideia geral, mas não para
publicar.
\item \textbf{O PSRF} (\emph{potential scale reduction factor}) \cite{gelman1992}. Na tabela do
\cmd{sump}, compara a variação de cada parâmetro dentro e entre as corridas. Deve ficar perto de
1,00 (na prática, até 1,01 ou 1,02).
\item \textbf{O ESS} (\emph{effective sample size}). As amostras de uma cadeia são
autocorrelacionadas: uma amostra se parece com a anterior. O ESS diz a quantas amostras
independentes as nossas equivalem. Deve ser maior que 200 (alguns aceitam 100). O \cmd{sump}
mostra o ESS mínimo (\cmd{minESS}) e o médio (\cmd{avgESS}) das duas corridas.
\end{enumerate}

\questao Quando a corrida terminar, procure no \arq{anolis_ND2_log.txt} o último valor do ASDSF,
e na tabela do \cmd{sump} o PSRF e o ESS de cada parâmetro. A corrida convergiu? Algum parâmetro
ficou mal amostrado?

\begin{resposta}
Na corrida de teste (100 mil gerações, 3 minutos): ASDSF final de 0,028; PSRF até 1,28
(\cmd{alpha\{1\}} e \cmd{pinvar\{1\}}); ESS mínimo de 10, com todos os 42 parâmetros abaixo de
200 e 39 abaixo de 100. Ou seja: não convergiu, como esperado para uma corrida tão curta. As corridas dos alunos vão variar um pouco, porque cada
uma começa de uma semente diferente. É uma boa deixa para a Questão~\ref{q:tracer}.
\end{resposta}

\paragraph{O Tracer.} Abra o Tracer e carregue os dois arquivos de parâmetros,
\arq{anolis_ND2.nex.run1.p} e \arq{anolis_ND2.nex.run2.p} (\ui{File > Import Trace File}). O
Tracer descarta 10\% como \emph{burn-in} por padrão; para ficar igual ao MrBayes, troque para
25\% clicando duas vezes no valor da coluna \ui{Burn-In} de cada arquivo e digitando 25000 (o
Tracer conta o \emph{burn-in} em gerações, e não em amostras). Para cada parâmetro, veja:
\begin{itemize}
\item na aba \ui{Trace}, o traço da cadeia ao longo das gerações: deve parecer uma ``taturana
peluda'', sem tendência de subir ou descer;
\item na aba \ui{Estimates}, o ESS (em vermelho quando é menor que 100, em amarelo entre 100 e
200);
\item selecionando as duas corridas juntas na aba \ui{Marginal Density}: as duas distribuições
devem se sobrepor.
\end{itemize}

\questao\label{q:tracer} No Tracer, olhe o traço do \cmd{LnL} e de dois parâmetros com ESS
baixo. Como eles se parecem? O que você faria para melhorar a amostragem?

\begin{resposta}
O caminho é rodar mais gerações. Com 1 milhão (32 minutos), o ASDSF caiu para 0,012, o ESS
mínimo subiu para 130 e o PSRF máximo para 1,012: quase lá. Com 3 milhões (a corrida que gerou a
árvore do mPTP, 1~h~36~min), ASDSF de 0,007, ESS mínimo de 416 e PSRF máximo de 1,002: todos os
critérios atendidos. Outra opção é simplificar o modelo das partições com parâmetros mal
amostrados.
\end{resposta}

\begin{quote}\small
\textbf{Na vida real.} As 100 mil gerações daqui são só para a corrida caber na aula. Numa
análise de verdade, rodaríamos pelo menos 10 milhões de gerações, e até todos os critérios serem
atendidos, num computador mais rápido ou num servidor como o CIPRES
(\url{https://www.phylo.org}). Também escolheríamos o esquema de partições e os modelos com um
programa como o PartitionFinder ou o ModelFinder, em vez de partir direto para as posições do
códon.
\end{quote}

\subsection{A árvore de consenso}\label{sec:consenso}

O MrBayes amostrou cerca de 2.000 árvores (1.000 por corrida), das quais cerca de 1.500 ficam depois do
\emph{burn-in}. Para resumir essas árvores numa só, o \cmd{sumt} faz uma \textbf{árvore de
consenso de maioria de 50\%}: entram na árvore os clados que aparecem em mais da metade das
árvores amostradas, e o número em cada nó é a probabilidade posterior do clado. Onde nenhum
arranjo passa de 50\%, a árvore fica com uma politomia. \textbf{É essa árvore, o consenso de
50\%, que normalmente interpretamos e publicamos numa análise com o MrBayes}, porque ela só
mostra os clados com mais apoio do que o contrário. Ela está no arquivo
\arq{anolis_ND2.nex.con.tre}.

O segundo \cmd{sumt}, com \cmd{contype=allcompat}, faz um consenso diferente: depois dos clados
com mais de 50\%, ele acrescenta, em ordem de frequência, os clados com menos de 50\% que ainda
forem compatíveis com os que já estão na árvore. O resultado é uma árvore quase totalmente
resolvida. Ela \emph{não} é para interpretar, porque mostra clados com pouco apoio como se
fossem tão confiáveis quanto os outros; só a fazemos porque o mPTP exige uma árvore sem
politomias (Seção~\ref{sec:preparar}). Ela está em \arq{anolis_ND2_allcompat.con.tre}.

Os comprimentos dos ramos nas duas árvores são as medianas do posterior, em substituições por
sítio.

\paragraph{O FigTree.} Abra o \arq{anolis_ND2.nex.con.tre} no FigTree. Quando ele perguntar
como chamar os rótulos dos nós, aceite o padrão. Depois:
\begin{itemize}
\item selecione o ramo de \cmd{BR\_ACB.CHUNB14548} e clique em \ui{Reroot} para enraizar no
grupo externo;
\item em \ui{Tree > Increasing Node Order}, ordene os nós para a árvore ficar mais fácil de ler;
\item marque \ui{Node Labels} no painel da esquerda e escolha \ui{prob} em \ui{Display}: é a
probabilidade posterior de cada clado (de 0 a 1).
\end{itemize}

\questao As sequências de \emph{A. chrysolepis} (\cmd{CR\_}) e de \emph{A. tandai} (\cmd{TA\_})
formam dois clados? Quantos clados bem apoiados (probabilidade posterior $\geq 0{,}95$) você vê
dentro de cada espécie? Há politomias? Onde?

\begin{resposta}
Nenhuma das duas espécies é monofilética nas árvores de teste, mas o arranjo muda de corrida
para corrida, e sempre com pouco apoio. São quatro grupos bem apoiados (probabilidade de 0,98
a 1,0): duas sequências de \emph{A. chrysolepis} (\cmd{CR\_SP.MRT6281} e \cmd{CR\_SP.MTR13939}),
as outras 56 de \emph{A. chrysolepis}, 46 de \emph{A. tandai} e 4 de \emph{A. tandai}. Na
corrida de 100 mil gerações, as 4 de \emph{A. tandai} ficam com as 56 de \emph{A.
chrysolepis} (0,54); nas de 1 e 3 milhões, as duas \cmd{CR\_SP} ficam com \emph{A. tandai}
(0,64 a 0,69). Detalhes no gabarito. Vale ligar isso à convergência: com pouca diferença entre
os arranjos, uma corrida curta não tem como decidir. Vale também chamar
a atenção para os ramos longos que separam grupos de sequências dentro de cada espécie: são
eles que o mPTP vai usar.
\end{resposta}

% =====================================================================
\section{Parte 2: delimitando linhagens com o mPTP}
% =====================================================================

\subsection{Preparando a árvore}\label{sec:preparar}

O mPTP precisa de uma árvore binária, enraizada e em formato Newick. O script
\arq{preparar_mptp.R} faz isso a partir do consenso \cmd{allcompat}:
\begin{enumerate}
\item lê a árvore (\cmd{read.nexus}, do pacote \cmd{ape});
\item enraíza em \emph{A. brasiliensis} (\cmd{root});
\item tira os dois grupos externos (\cmd{drop.tip}): eles são espécies diferentes, distantes
das outras, e não queremos delimitá-los; o mPTP até tem opções para isso
(\cmd{-{}-outgroup} e \cmd{-{}-outgroup\_crop}), mas tirá-los antes é mais simples;
\item se ainda houver alguma politomia, resolve-a ao acaso, com ramos novos de comprimento
zero (\cmd{multi2di}). Na árvore do professor não há nenhuma, e o passo não muda nada. Mas o
\cmd{allcompat} de uma corrida curta, como a de vocês, pode sair com politomias (quando os
clados que as resolveriam nunca foram amostrados), e aí o mPTP se recusa a ler a árvore. Os
ramos novos têm comprimento zero, menor que o ramo mínimo (próxima seção), e o mPTP não os
conta; mas eles decidem quais sequências podem formar uma linhagem juntas, e por isso o script
fixa a semente do sorteio;
\item grava a árvore em Newick (\arq{anolis_mptp.tre}) e um FASTA sem os grupos externos
(\arq{anolis_mptp.fas}), que o mPTP vai usar para calcular o ramo mínimo.
\end{enumerate}

Para todos chegarem ao mesmo resultado, o professor já preparou esses dois arquivos a partir de
uma corrida de 3 milhões de gerações, e é com eles que vamos trabalhar. Se você quiser preparar a
sua própria árvore, ponha o \arq{preparar_mptp.R} na pasta da corrida, abra no RStudio, use
\ui{Session > Set Working Directory > To Source File Location} e rode. O script vai gravar por
cima dos arquivos do professor; guarde uma cópia deles antes.

\begin{resposta}
O \cmd{allcompat} das corridas de 100 mil e de 3 milhões de gerações já saiu totalmente
resolvido (108 nós internos, 38 e 39 deles com probabilidade posterior abaixo de 0,5), e o
\cmd{multi2di} não mudou nada. Mantive o passo porque, numa corrida de teste de 20 mil
gerações, o \cmd{allcompat} saiu com uma politomia de quatro ramos, e o mPTP não leu a árvore.
\end{resposta}

\subsection{Instalando o mPTP}\label{sec:instalarmptp}

O mPTP é um programa de linha de comando, sem janela: tudo é feito digitando comandos.
\begin{itemize}
\item \textbf{Windows.} Baixe o arquivo para Windows em
\url{https://github.com/Pas-Kapli/mptp/releases}, descompacte, e copie o \arq{mptp.exe} para a
pasta da prática. O mPTP do Windows roda pelo \textbf{PowerShell}: abra a pasta da prática no
Explorador de Arquivos, clique com o botão direito num espaço vazio e escolha
\ui{Abrir no Terminal} (ou, no Windows 10, segure Shift e escolha
\ui{Abrir janela do PowerShell aqui}). Nos comandos abaixo, troque \cmd{mptp} por
\cmd{.\textbackslash mptp.exe}; por exemplo, \cmd{.\textbackslash mptp.exe -{}-help}.
\item \textbf{Linux.} Baixe o arquivo \arq{.tar.gz} para Linux na mesma página, descompacte, e
o programa é o \arq{mptp} da pasta \arq{bin}.
\item \textbf{WSL.} Quem fizer a prática no WSL segue as instruções do Linux. Para instalar o
WSL: no PowerShell como administrador, \cmd{wsl -{}-install}, reinicie e abra o Ubuntu. As pastas
do Windows ficam em \arq{/mnt/c/} (por exemplo, \arq{/mnt/c/Users/seu_nome/Downloads}).
\item \textbf{Mac.} Baixe o \arq{mptp-0.2.5-macos-aarch64.tar.gz} na mesma página (é a versão
para os Macs com processador Apple, M1 ou mais novo), descompacte com dois cliques, e copie o
\arq{mptp} da pasta \arq{bin} para a pasta da prática. No Terminal, dentro da pasta, rode uma
vez:
\begin{Verbatim}
xattr -d com.apple.quarantine mptp
chmod +x mptp
\end{Verbatim}
O primeiro comando tira o bloqueio que o Mac põe em programas baixados; o segundo deixa o
arquivo executável. Depois, chame o programa como \cmd{./mptp}.
\end{itemize}
Teste com \cmd{mptp -{}-help} (\cmd{./mptp -{}-help} no Mac e no Linux, se o programa estiver
na pasta da prática; \cmd{.\textbackslash mptp.exe -{}-help} no Windows). Os comandos abaixo
foram quebrados em mais de uma linha só para caber na página: no Terminal ou no PowerShell,
digite cada um numa linha só.

\begin{resposta}
O executável do Mac é só para Apple Silicon (\cmd{arm64}). Se aparecer algum Mac Intel
(confira em \ui{Sobre este Mac}), dá para compilar o mPTP nele (precisa do \cmd{flex} e do
\cmd{bison}: \cmd{brew install flex bison}) ou passar o seu executável compilado. No Windows,
se o PowerShell bloquear o programa, clique com o botão direito no \arq{mptp.exe},
\ui{Propriedades}, e marque \ui{Desbloquear}. Os resultados do gabarito foram feitos com o
mPTP compilado no Linux; com a mesma semente, os executáveis do Windows e do Mac devem dar a
mesma delimitação, mas os números do MCMC podem variar um pouco.
\end{resposta}

\subsection{O comprimento mínimo de ramo}

Sequências idênticas ou quase idênticas ficam ligadas por ramos muito curtos, que não refletem
eventos de coalescência e atrapalham o modelo. O mPTP ignora os ramos mais curtos que um valor
mínimo, e esse valor pode ser calculado a partir do alinhamento:
\begin{Verbatim}
mptp --minbr_auto anolis_mptp.fas --tree_file anolis_mptp.tre
     --output_file anolis_minbr
\end{Verbatim}
(tudo numa linha só). O valor aparece na tela, numa linha que começa com \cmd{Minimum branch
length (-{}-minbr) should be set to}, e no arquivo \arq{anolis_minbr.txt}. Use o valor do
arquivo, com seis casas decimais, para todos chegarem aos mesmos números.

\questao Qual é o comprimento mínimo de ramo sugerido? A quantas substituições no ND2 inteiro
(924 sítios) ele corresponde, aproximadamente?

\begin{resposta}
0,001587. Vezes 924 sítios, dá cerca de 1,5 substituição: ramos com menos do que isso (na
prática, sem nenhuma ou com uma só substituição) são tratados como ``nada aconteceu''.
\end{resposta}

\subsection{A delimitação de máxima verossimilhança}

Com o valor do ramo mínimo (troque \cmd{X} pelo valor da seção anterior):
\begin{Verbatim}
mptp --ml --multi --tree_file anolis_mptp.tre --minbr X
     --output_file anolis_ml
\end{Verbatim}
Na tela, e no arquivo \arq{anolis_ml.txt}:
\begin{itemize}
\item \cmd{Null-model score} e \cmd{Best score for multi coalescent rate}: os logaritmos da
verossimilhança do modelo nulo (todas as sequências são uma espécie só) e da melhor
delimitação;
\item \cmd{LRT computed p-value}: o teste da razão de verossimilhanças entre os dois; se
\cmd{LRT: passed}, a delimitação é significativamente melhor que o modelo nulo;
\item \cmd{Number of delimited species}: o número de linhagens;
\item depois, os blocos \cmd{Species 1:}, \cmd{Species 2:}, \ldots, com as sequências de cada
linhagem.
\end{itemize}
O arquivo \arq{anolis_ml.svg} mostra a árvore com as linhagens: abra no navegador (Chrome,
Firefox, Safari). Nele, os ramos do processo entre espécies aparecem em verde, e os ramos de
dentro de cada linhagem, em vermelho: cada grupo vermelho é uma linhagem.

\questao Quantas linhagens o mPTP encontrou? Alguma linhagem mistura sequências de
\emph{A. chrysolepis} e de \emph{A. tandai}? Alguma tem uma sequência só?

\begin{resposta}
15 linhagens (LRT: passed), 7 só com \emph{A. tandai} e 8 só com \emph{A. chrysolepis};
nenhuma mistura as duas espécies (as duas \cmd{CR\_SP} que caem junto de \emph{A. tandai} na
árvore formam uma linhagem própria). Duas linhagens têm uma sequência só: \cmd{TA\_KX760654} e
\cmd{CR\_HT5068}. A lista completa está no gabarito.
\end{resposta}

\subsection{O suporte da delimitação}

A delimitação de máxima verossimilhança é uma só, mas outras delimitações podem ser quase tão
boas. Para saber quanto confiar em cada linhagem, o mPTP faz um MCMC sobre as delimitações
possíveis, nesta mesma árvore. Vamos rodar duas cadeias, de 10 milhões de gerações cada,
amostrando a cada 10.000 e descartando o primeiro milhão:
\begin{Verbatim}
mptp --mcmc 10000000 --multi --mcmc_runs 2 --mcmc_sample 10000
     --mcmc_burnin 1000000 --tree_file anolis_mptp.tre --minbr X
     --output_file anolis_mcmc --seed 2026
\end{Verbatim}
Leva poucos segundos. O MCMC do mPTP é bem mais simples que o do MrBayes, porque a árvore está
fixa: só a delimitação muda. Na tela, no fim:
\begin{itemize}
\item \cmd{CCI} e \cmd{HPD}: o intervalo de credibilidade do número de linhagens;
\item \cmd{ML average support based on run with seed ...}: o suporte médio das linhagens da
delimitação de máxima verossimilhança, em cada corrida;
\item \cmd{Average standard deviation of support values among runs}: a diferença entre as duas
corridas (deve ser perto de zero, como o ASDSF do MrBayes).
\end{itemize}
O arquivo \arq{anolis_mcmc.2026.combined.svg} mostra a árvore com o suporte das duas corridas
juntas: a cor vai do preto (suporte 1, nó entre linhagens) ao vermelho (suporte 0, nó dentro de
uma linhagem), e o valor aparece escrito nos nós com suporte acima de 0,5. Os valores de todos
os nós estão no \arq{anolis_mcmc.2026.combined.tree}, que abre no FigTree (escolha
\ui{label} em \ui{Node Labels}). O suporte de cada nó é a fração das amostras do MCMC em que aquele nó estava no
processo \emph{entre} espécies, ou seja, a fração em que ele separava linhagens diferentes. Para
cada linhagem, olhe dois nós:
\begin{itemize}
\item o nó logo \emph{acima} da linhagem (o ancestral dela com a linhagem vizinha): perto de 1
quer dizer que ela foi separada das vizinhas em quase todas as amostras;
\item o nó da \emph{base} da linhagem (o ancestral comum das sequências dela): perto de 0 quer
dizer que ela ficou inteira em quase todas as amostras; um valor alto quer dizer que, em muitas
amostras, o mPTP a dividiu em mais de uma.
\end{itemize}
A mesma informação, resumida para o número de linhagens, está nos arquivos \arq{.stats}: a
porcentagem das amostras com 14, 15, 16\ldots\ linhagens.

\questao Qual é o suporte médio da delimitação de máxima verossimilhança? Pelo arquivo
\arq{.stats}, o número de linhagens da delimitação de máxima verossimilhança é também o mais
frequente no MCMC? Quais linhagens o MCMC costuma dividir?

\begin{resposta}
Suporte médio de 0,91 nas duas corridas. As 15 linhagens da máxima verossimilhança aparecem em
só 2,5\% das amostras; o mais frequente são 16 (27\%) e 18 (20\%), e o intervalo vai de 14 a 22.
Todas as linhagens estão bem separadas das vizinhas (nó acima com suporte de 0,999 ou 1), e a
incerteza está dentro de duas linhagens grandes de \emph{A. tandai}, que o MCMC divide com
frequência (nó da base com 0,63 e 0,41). É uma boa deixa para lembrar que a delimitação de
máxima verossimilhança é um ponto, e o MCMC mostra a incerteza em volta dele.
\end{resposta}

\begin{quote}\small
\textbf{Na vida real.} O resultado depende da árvore: valeria repetir o mPTP com a árvore de
máxima verossimilhança (de um IQ-TREE ou RAxML) e com árvores amostradas do posterior, e
comparar. Também valeria comparar com outros métodos de delimitação (ABGD, ASAP, GMYC), e,
sobretudo, testar as linhagens com genes nucleares e outros dados antes de chamar qualquer uma
de espécie.
\end{quote}

% =====================================================================
\appendix
\section{Alterações em relação ao tutorial original}\label{ap:mudancas}
% =====================================================================

\begin{enumerate}
\item \textbf{Dados.} No lugar da matriz de Hylodidae (24 terminais, vários genes), o ND2 de 110
\emph{Anolis}, com mais sequências por espécie, como é preciso para a delimitação.
\item \textbf{Partições e modelos.} Três partições pelas posições do códon, com
\cmd{nst=mixed} (o MrBayes amostra o modelo de cada partição), em vez de modelos escolhidos
antes no PartitionFinder.
\item \textbf{Corrida.} Só local, sem a etapa do CIPRES, e com o MrBayes 3.2.9.
\item \textbf{Convergência.} Acrescentado o Tracer para o ESS e os traços.
\item \textbf{Consenso.} Explicados o consenso de maioria de 50\% e o \cmd{allcompat}, que só
serve para o mPTP.
\item \textbf{Nova parte.} Acrescentada a delimitação com o mPTP.
\end{enumerate}

% =====================================================================
\begin{thebibliography}{99}
\bibitem{gelman1992} Gelman, A. \& Rubin, D.~B. 1992. Inference from iterative simulation
using multiple sequences. \emph{Statistical Science} 7: 457--472.
\bibitem{huelsenbeck2004} Huelsenbeck, J.~P., Larget, B. \& Alfaro, M.~E. 2004. Bayesian
phylogenetic model selection using reversible jump Markov chain Monte Carlo. \emph{Molecular
Biology and Evolution} 21: 1123--1133.
\bibitem{kapli2017} Kapli, P., Lutteropp, S., Zhang, J., Kobert, K., Pavlidis, P., Stamatakis, A.
\& Flouri, T. 2017. Multi-rate Poisson tree processes for single-locus species delimitation
under maximum likelihood and Markov chain Monte Carlo. \emph{Bioinformatics} 33: 1630--1638.
\bibitem{oliveira2025} Vieira de Oliveira, A.~P., D'Angiolella, A., Azevedo, J.~A.~R., Vitt,
L.~J., Colli, G.~R., Choueri, E.~L., Gamble, T., Taucce, P.~P.~G., Rodrigues, M.~T., Noonan,
B.~P., Avila-Pires, T.~C. \& Werneck, F.~P. 2025. The influence of historical climate cycles on
the diversification and demography of species from the Neotropical lizard group \emph{Anolis
chrysolepis}. \emph{Biological Journal of the Linnean Society} 144: blaf006.
\bibitem{paradis2019} Paradis, E. \& Schliep, K. 2019. ape 5.0: an environment for modern
phylogenetics and evolutionary analyses in R. \emph{Bioinformatics} 35: 526--528.
\bibitem{rambaut2018} Rambaut, A., Drummond, A.~J., Xie, D., Baele, G. \& Suchard, M.~A. 2018.
Posterior summarization in Bayesian phylogenetics using Tracer 1.7. \emph{Systematic Biology}
67: 901--904.
\bibitem{ronquist2012} Ronquist, F., Teslenko, M., van der Mark, P., Ayres, D.~L., Darling, A.,
Höhna, S., Larget, B., Liu, L., Suchard, M.~A. \& Huelsenbeck, J.~P. 2012. MrBayes 3.2:
efficient Bayesian phylogenetic inference and model choice across a large model space.
\emph{Systematic Biology} 61: 539--542.
\bibitem{zhang2013} Zhang, J., Kapli, P., Pavlidis, P. \& Stamatakis, A. 2013. A general
species delimitation method with applications to phylogenetic placements.
\emph{Bioinformatics} 29: 2869--2876.
\end{thebibliography}

\end{document}
